\section{Classification of the specified connected networks}
The three-cell computation has an exact extension to every larger connected graph.
Network-controllability arguments and algebraic propagation are established tools~\cite{Burgarth2009}; the point here is to close the extension for the two specified bonds and the restricted local algebra.

\begin{lemma}[Extension of a fully controlled block]\label{lem:extension}
Suppose a block of $m\geq2$ cells generates $\mathfrak{su}(d)$, where $d=4^m$.
Append a four-level cell with its independent local $\mathfrak{spin}(5)$ controls.
If a connecting control is $A_0\otimes\Gamma_5$ or $A_0\otimes\Sigma_{45}$ with $A_0$ a nonzero traceless Hermitian block operator, the enlarged control algebra is $\mathfrak{su}(4d)$.
\end{lemma}
\begin{proof}
Write $V=\operatorname{span}_{\mathbb R}\{\Gamma_i\}$ and $B=\operatorname{span}_{\mathbb R}\{\Sigma_{ij}\}$.
The real adjoint module of the compact simple algebra $\mathfrak{su}(d)$ is irreducible.
The local $V$ and $B$ modules are respectively the irreducible vector and adjoint modules of $\mathfrak{spin}(5)$.
Independent repeated commutation with block and local controls therefore generates all of $\mathfrak{su}(d)\otimes V$ from a vector bond, or all of $\mathfrak{su}(d)\otimes B$ from a bivector bond.

In the vector case choose two commuting block Pauli involutions $A,B_0$ whose product is nonzero and traceless, such as two independent $Z$ strings.
Then
\begin{equation}
 \Hbr{A\otimes\Gamma_1}{B_0\otimes\Gamma_2}
 =2AB_0\otimes\Sigma_{12}.
\end{equation}
Taking the block adjoint span and the local bivector orbit gives $\mathfrak{su}(d)\otimes B$ too.
Conversely, in the bivector case choose anticommuting block Pauli involutions $A,B_0$ and the commuting local bivectors $\Sigma_{12},\Sigma_{34}$.
The Clifford volume convention gives $\Sigma_{12}\Sigma_{34}=-\Gamma_5$, hence
\begin{equation}
 \Hbr{A\otimes\Sigma_{12}}{B_0\otimes\Sigma_{34}}
 =-\Hbr{A}{B_0}\otimes\Gamma_5.
\end{equation}
This generates the vector module as well.

Choose a traceless block involution $Q$ with $Q^2=I$.
Since both tensor modules are available,
\begin{equation}
 \Hbr{Q\otimes\Sigma_{12}}{Q\otimes\Gamma_1}
 =2I\otimes\Gamma_2.
\end{equation}
Local rotations supply every $I\otimes\Gamma_i$.
Together with the original $I\otimes B$, these form $I\otimes\mathfrak{su}(4)$, because the five vectors and ten bivectors span all traceless Hermitian $4\times4$ matrices.
Independent block and full local adjoint actions on any nonzero tensor then supply $\mathfrak{su}(d)\otimes\mathfrak{su}(4)$.
The three resulting subspaces have total dimension
\begin{equation}
 (d^2-1)+15+15(d^2-1)=16d^2-1,
\end{equation}
which exhausts $\mathfrak{su}(4d)$.
\end{proof}

\begin{theorem}[Connected-network control algebra]\label{thm:graph}
On a finite connected graph, allow independent copies of the ten local generators and independently available nonzero edge controls of type $S_rS_s$ or $C_rC_s$.
One cell generates $\mathfrak{spin}(5)$.
Two cells with only an $S$ bond generate $\mathfrak{spin}(10)$; with a $C$ bond present they generate $\mathfrak{su}(16)$.
Every connected graph of $n\geq3$ cells with these edge types generates $\mathfrak{su}(4^n)$.
\end{theorem}
\begin{proof}
The one-, two- and three-cell bases are the exact Pauli closures in Theorem~\ref{c:thm:proliferation}.
If a $C$ edge is present, its two endpoints already give a fully controlled $\mathfrak{su}(16)$ block.
If every edge is of type $S$, a connected graph with at least three vertices contains a three-vertex path, whose controls generate $\mathfrak{su}(64)$.
In either case grow the controlled block along a spanning tree.
At each step the block part of the next edge is a nonzero traceless operator, so Lemma~\ref{lem:extension} applies.
Extra available edges cannot reduce a generated algebra.
\end{proof}

The assumption of independently available controls matters.
A single fixed homogeneous Hamiltonian, a diagonal collective control restriction or a projected low-energy model is a different problem.
The theorem rules out the naive repetition of this pair as a restricted spatial control algebra; it does not rule out every constrained or collective model.
As a portability check, the integer algorithm also gives $65535=\dim\mathfrak{su}(256)$ for four-cell chains with either uniform bond.
This check supports the implementation; the analytic induction, rather than an extrapolation of dimensions, proves the all-graph result.
